\ResumenAbstract
% RESUMEN EN ESPAÑOL
{
En el contexto de la rápida evolución tecnológica, la Educación Superior enfrenta el
desafío de diseñar Planes de Estudio flexibles que respondan a la demanda de
perfiles profesionales especializados. Esta investigación aborda este problema
aplicando los principios de la Ingeniería de Líneas de Productos de Software (SPL) y
los Modelos de Características (Feature Models - FM) al dominio de la planificación
curricular.\\
El trabajo se centra en el diseño, implementación y validación de un \textit{backend} que
permita a los gestores académicos modelar currículos como "Líneas de Productos
Educativos". El sistema implementa una API RESTful que sirve como núcleo para las
operaciones de creación, análisis, validación y configuración de FMs. La arquitectura
del servidor gestiona la lógica de negocio para representar
características y relaciones de obligatoriedad, optatividad y prerrequisitos como
restricciones formales, garantizando la consistencia e integridad de los datos
mediante un motor de reglas eficiente y un modelo de datos optimizado.\\
El resultado principal es un servicio backend especializado que no solo soporta el
modelado y la persistencia de FMs adaptados al ámbito educativo, sino que también
proporciona la base computacional para automatizar la generación y validación de
Planes de Estudio diversos y consistentes a partir de un tronco común. Este
desarrollo sienta las bases técnicas para una nueva generación de herramientas de
gestión académica, demostrando la viabilidad de transferir técnicas avanzadas de
ingeniería de software al dominio educativo.\\
\textbf{Palabras claves:} Modelos Característicos, Líneas de Productos de
Software, Líneas de Productos Educativos, Variabilidad, Motor de Reglas, Plan de Estudios.
}
% RESUMEN EN INGLÉS
{
In the context of rapid technological evolution, Higher Education faces the challenge of designing flexible Study Plans that respond to the demand for specialized professional profiles. This research addresses this problem by applying the principles of Software Product Line Engineering (SPL) and Feature Models (FM) to the domain of curriculum planning.\\
The work focuses on the design, implementation, and validation of a backend that allows academic managers to model curricula as "Educational Product Lines". The system implements a RESTful API that serves as the core for the creation, analysis, validation, and configuration of FMs. The server architecture manages the business logic to represent subjects as features and relationships of obligation, optionality, and prerequisites as formal constraints, ensuring data consistency and integrity through an efficient rule engine and an optimized data model.\\
The main result is a specialized backend service that not only supports the modeling and persistence of FMs adapted to the educational field but also provides the computational basis for automating the generation and validation of diverse and consistent Study Plans from a common trunk. This development lays the technical foundations for a new generation of academic management tools, demonstrating the feasibility of transferring advanced software engineering techniques to the educational domain.\\
\textbf{Keywords:} Feature Models, Software Product Lines, Educational Product Lines, Variability, Rule Engine, Curriculum.
}
